\documentclass[11pt,a4paper]{article}  % Déclare la classe du document.

\newcommand{\chaptermark}[2]{\og #2 \fg{} (#1)} % ligne uniquement pour éviter une erreur en cas d'article
\include{../1ere-preambule}

\usepackage{epic}   % extension et simplification de  l'environnement picture
\usepackage{graphicx}    % permet d'inclure des graphique externe dans un document
\usepackage{arydshln}   %permet d'avoir des lignes en pointillés

\newcommand{\lignes}[1]{
	\setlength{\unitlength}{1mm}
	\vskip 0,4cm\noindent
	\multido{\i=0+1}{#1}{%
		\noindent
		\begin{picture}(150,5)
			\linethickness{0,005mm}
			\put(-5,0){\dashline{1}(175,5)(5,5)}
		\end{picture}\vskip 0,001cm
	}
	\vskip 0cm
}

\rhead{Accompagnement - Séance 1 - Equations}

\newcommand{\va}[1]{\lvert \ #1 \ \rvert}
\newcommand{\vaf}[1]{\left\lvert \ #1\  \right\rvert}

\begin{document}
	\graphicspath{{images/}}
	
	\begin{center}
		
		
		\fbox{
			\begin{minipage}[]{15cm}
				\begin{center}
					\LARGE{\rule{0.5cm}{0pt}\rule[-0.25cm]{0pt}{0.9cm}
						Accompagnement - Séance 1 - 03/10/2022\\
						Equations et inéquations
					}
				\end{center}
			\end{minipage}
		}
	\end{center}

\exercice{ - équations du 1er degré}\\
Résoudre dans \R les équations.
\begin{multicols}{2}
	\begin{enumMet}
		\item $4 x-15=x+1$ ;
		\item $-5 x-2=-x+18$ ;
		\item $1,2 x+0,3=0,4$ ;
		\item $\dfrac{2}{3} x+\dfrac{1}{5}=\dfrac{3}{4} x+\dfrac{1}{3}$;
	\end{enumMet}
\end{multicols}

\exercice{ - équations du 2nd degré}\\
Résoudre dans \R les équations.
\begin{multicols}{2}
	\begin{enumMet}
		\item $(7 x+2)(x+14)=0$ ;
		\item $(6-5 x)(4 x+1)=0$ ;
		\item $x(4 x+1)+7(4 x+1)=0$ ;
		\item $5(x-3)^3=7 x(x-3)^2$;
		\item $(x+3)^2=(2 x-5)^2$;
		\item $\left(x^2+1\right)^2=\left(3 x^2-1\right)^2$;
		\item $4(1-x)(4 x+9)(2 x+3)=0$;
		\item $10(x+7)(x-5)=3 x(x+7)$
	\end{enumMet}
\end{multicols}

\exercice{ - équations quotient}\\
Résoudre dans \R les équations.
\begin{multicols}{2}
	\begin{enumMet}
		\item $\dfrac{2 x+1}{x-1}=\dfrac{3}{5}$ ;
		\item $\dfrac{x^2-3 x+9}{3 x^2+5}=\dfrac{1}{3}$ ;
		\item $\dfrac{4 x-1}{\left(x^2+x+1\right)^2}=0$ ;
		\item $\dfrac{x}{x-1}=\dfrac{2 x+3}{2 x}$;
		\item $\dfrac{3 x+2}{x+4}=2$;
		\item $\dfrac{1}{x}+\dfrac{1}{x-4}=0$;
	\end{enumMet}
\end{multicols}

\exercice{ - résolution d'inéquations du premier degré}\\
Résoudre dans $\R$  les inéquations suivantes.
\begin{multicols}{2}
	\begin{enumMet}
		\item $3(2-x)<2(7-3 x)$ ;
		\item $\dfrac{3}{2} x+2 \geqslant-2 x+\dfrac{4}{3}$ ;
		\item $\dfrac{x+3}{2} \geqslant \dfrac{x+1}{3}$ ;
		\item $-\dfrac{3}{8} x+2>\dfrac{1}{4} x-1$;
	\end{enumMet}
\end{multicols}
\vspace*{-0.3cm}

\exercice{ - résolution d'inéquations du second degré}\\
Résoudre dans $\R$  les équations suivantes.
\begin{multicols}{2}
    \begin{enumMet}
        \item $x^2(x+4) \leqslant 0$ ;
        \item $\left(5+x^2\right)(x-7) \geqslant 0$ ;
        \item $\left(3-\dfrac{x}{2}\right)\left(x+\dfrac{20}{3}\right) \leqslant 0$ ;
        \item $(x+12)(1-x)<-(x+12)(3 x+7)$;
        \item $(3 x+5)^2 \leqslant(2 x+1)^2$;
        \item $\dfrac{x}{x-16} \leqslant 0$
        \item $\dfrac{x-3}{x+3} \geqslant 0$
        \item $1-\dfrac{1}{x}<4$
        \item $\dfrac{x}{6 x-1} \leqslant \dfrac{5}{3}$
    \end{enumMet}
\end{multicols}

\end{document}

